\scanpage{20}\mspage{- 16 -}

\uline{Preuve} : (I) D'après (6.3), il n'y a en effet pas plus de deux automorphismes.

(II) et (III) Lorsque les courbes sont mises sous la forme (2.1), il résulte de (2.5) que les automorphismes sont décrits parleur action sur $\pi$ et que $\lambda\in G_m$ défini un automorphisme si et seulement si
\[
c_4=\lambda^4c_4\qquad\text{et}\qquad c_6=\lambda^6c_6\quad,
\]
d'où les assertions.

(IV) Si la courbe est mise sous forme (6.1), les automorphismes sont donnés par (1.2) : $x'\handwritten{\rightsquigarrow}x$, pour $s=t=0$ et
\[
\handwritten{\Bigg\{}\begin{aligned}
&u^4=1\\
&r^3+ra_4+a_6(1-u^2)=0
\end{aligned}
\]
et la loi de composition par
\[
(u',r')\quad(u,r)=(u\,u',\,u^{\prime 2}r+r')
\]

Lorsque la courbe est mise sous forme (6.7), ces équations se réduisent à $u\in\mu_4$ et $r\in F_3$.

(V) Si la courbe est mise sous forme (6.3), les automorphismes sont donnés par (1.2) : $x'\handwritten{\rightsquigarrow}x$, pour $r=s^2$ et
\[
\handwritten{\Bigg\{}\begin{aligned}
&u^3=1\\
&s^4+sa_3+a_4(1-u)=0\\
&t^2+ta_3+s^6+s^2a_4=0
\end{aligned}
\]

% archive sheet number: 19
